\chapter{Game Design Principles and the \prg Framework}
\label{ch:platform}

This is the one chapter that belongs to both tracks. Whether you are here for
assembly or for C, you need a shared mental model of two things: how a small game
is structured as a program, and how the \prg framework turns that structure into
something you can see and hear. Read it once now; you will return to it often.

\section{What a game actually is}

Strip away the graphics and a game is a very old idea: a loop that never ends
until you tell it to. Each turn around the loop is called a \emph{frame}. In a
single frame the machine does three things, always in the same order. It
\emph{reads} the state of the world---which buttons are pressed, how much time has
passed. It \emph{updates} an internal model of the game---moves the ball, checks
for collisions, adjusts the score. And it \emph{draws} the new state to the
screen. Then it does it again, fast enough that the eye perceives motion rather
than a slideshow.

\begin{prgconcept}
A game is a loop over frames. In each frame the program reads input, updates its
state, and draws the result. ``State'' is simply the set of variables that
describe the game right now: positions, velocities, scores, lives. Everything
interesting a game does is a rule for turning this frame's state into next
frame's state.
\end{prgconcept}

It is worth dwelling on the discipline this imposes, because it is the discipline
\prg enforces. Update and draw are kept separate on purpose. Update is allowed to
change state but should not draw; draw is allowed to read state and put pixels on
the screen but should not change the game's logic. A beginner who mixes the two
ends up with games where, for instance, a collision is detected only when an
object happens to be drawn, and bugs of that kind are maddening to find. Keeping
the phases separate is not bureaucracy; it is the single most effective habit for
writing games that behave predictably.

\begin{center}
\begin{tikzpicture}[font=\sffamily\small,node distance=6mm,
  phase/.style={rounded corners=3pt,draw=prgsteel,fill=prgpanel,minimum width=26mm,minimum height=11mm,align=center},
  once/.style={rounded corners=3pt,draw=prggreen,fill=prggreen!10,minimum width=26mm,minimum height=11mm,align=center},
  ar/.style={-{Stealth[length=2.2mm]},draw=prgink!70,thick}]
  \node[once] (init) {\textbf{init}\\\scriptsize once at start};
  \node[phase,right=14mm of init] (read) {\textbf{read input}\\\scriptsize buttons, time};
  \node[phase,right=10mm of read] (update) {\textbf{update}\\\scriptsize change state};
  \node[phase,right=10mm of update] (draw) {\textbf{draw}\\\scriptsize render state};
  \draw[ar] (init) -- (read);
  \draw[ar] (read) -- (update);
  \draw[ar] (update) -- (draw);
  \draw[ar] (draw.south) .. controls ([yshift=-9mm]draw.south) and ([yshift=-9mm]read.south) .. (read.south)
     node[midway,below,font=\scriptsize\itshape]{next frame};
\end{tikzpicture}
\captionof{figure}{The frame loop. \prg calls your \code{init} once, then your
\code{update} and \code{draw} once per frame, forever.}
\label{fig:frameloop}
\end{center}

\section{The three-function contract}

\prg expresses the frame loop as a contract of exactly three functions that you
provide and the framework calls \citep{prg32teaching}. Every example in this
book, in either language, exports the same three entry points, distinguished only
by a prefix that names the game:

\begin{lstlisting}[style=plain,caption={The universal shape of a \prg game. The prefix changes; the structure never does.}]
<prefix>_init     -- runs once when the game (cartridge) is selected
<prefix>_update   -- runs once per frame; changes game state
<prefix>_draw     -- runs once per frame; renders the current state
\end{lstlisting}

For a Pong game written in assembly the prefix is \code{pong\_graphics}, so the
symbols are \code{pong\_graphics\_init}, \code{pong\_graphics\_update}, and
\code{pong\_graphics\_draw}. For the C version the prefix is \code{pong\_c}. The
framework does not care which language produced the symbols; it only needs the
three of them.

\begin{prgconcept}
The entire interface between your game and \prg is three functions:
\code{init}, \code{update}, and \code{draw}. Learn this shape once and every
example in the book---assembly or C, emulator or hardware---becomes a variation
on it.
\end{prgconcept}

\section{The screen, the colours, and the coordinate system}

\prg presents a display that is 320 pixels wide and 240 pixels tall, but games
draw into a centred \emph{viewport} of 320 by 200, leaving thin bands at the top
and bottom for status information \citep{prg32repo}. The coordinate origin
$(0,0)$ is the top-left corner; $x$ increases to the right and $y$ increases
\emph{downward}. This downward-$y$ convention surprises newcomers who remember
graphs from mathematics, where $y$ goes up, and it is one of the most common
sources of early confusion.

\begin{prgwarn}
On the screen, $y$ grows \emph{downward}. A larger $y$ means lower on the
display. If your object moves up when you expected it to move down, you have very
likely confused the sign of a $y$ change. This is normal; check it first.
\end{prgwarn}

\begin{center}
\begin{tikzpicture}[font=\sffamily\scriptsize,scale=0.95]
  % full 320x240 frame
  \draw[draw=prgink!50,fill=prgink!4] (0,0) rectangle (8,6);
  % top and bottom bands
  \fill[prgsteel!25] (0,5.5) rectangle (8,6);
  \fill[prgsteel!25] (0,0) rectangle (8,0.5);
  % viewport 320x200 centered
  \draw[draw=prgaccent,thick,fill=prgaccent!6] (0,0.5) rectangle (8,5.5);
  % origin marker
  \fill[prgred] (0,5.5) circle (1.4pt);
  \node[anchor=south west,text=prgred] at (0.05,5.55){(0,0)};
  % axes
  \draw[-{Stealth[length=2mm]},prgink] (0.3,5.2) -- (2.3,5.2) node[right]{$x \to$};
  \draw[-{Stealth[length=2mm]},prgink] (0.3,5.2) -- (0.3,3.6) node[below right]{$y \downarrow$};
  \node[text=prgaccent] at (4,3){game viewport \textbf{320 $\times$ 200}};
  \node[text=prgsteel] at (4,5.75){top status band};
  \node[text=prgsteel] at (4,0.25){bottom status band};
  \node[anchor=north west] at (0,-0.15){\itshape Full display: 320 $\times$ 240};
\end{tikzpicture}
\captionof{figure}{The \prg display. Framework screens use the full
$320\times240$; your game draws into the centred $320\times200$ viewport.}
\label{fig:screen}
\end{center}

Colours are expressed in the \textbf{RGB565} format: a 16-bit number that packs
five bits of red, six of green, and five of blue. You rarely need to compute these
by hand, because the framework defines names for the common ones---\code{PRG32\_COLOR\_BLACK}
is \code{0x0000}, \code{PRG32\_COLOR\_WHITE} is \code{0xffff},
\code{PRG32\_COLOR\_YELLOW} is \code{0xffe0}, and so on \citep{prg32repo}. When you
see the literal \code{65535} in an assembly example, it is simply white written
in decimal.

\begin{prgaside}
RGB565 was the natural pixel format of countless devices in the era this book
celebrates, precisely because it fits a pixel in two bytes while still giving the
eye enough colour. The extra green bit reflects the fact that human vision is most
sensitive to green. The format is a small, elegant compromise between memory and
fidelity---exactly the kind of constraint that makes retro programming
instructive.
\end{prgaside}

\section{Reading the player: the input bitmask}

\prg reports all the buttons at once as a single integer, a \emph{bitmask}, in
which each button is one bit \citep{prg32repo}. Reading input is a single call,
\code{prg32\_input\_read}, that returns this integer. To find out whether a
particular button is pressed, you test its bit with a bitwise AND.

\begin{center}
\begin{tabular}{@{}llr@{}}
\toprule
Constant & Button & Bit value \\
\midrule
\code{PRG32\_BTN\_LEFT}   & left  & \code{0x01} \\
\code{PRG32\_BTN\_RIGHT}  & right & \code{0x02} \\
\code{PRG32\_BTN\_UP}     & up    & \code{0x04} \\
\code{PRG32\_BTN\_DOWN}   & down  & \code{0x08} \\
\code{PRG32\_BTN\_A}      & A     & \code{0x10} \\
\code{PRG32\_BTN\_B}      & B     & \code{0x20} \\
\code{PRG32\_BTN\_SELECT} & select & \code{0x40} \\
\bottomrule
\end{tabular}
\captionof{table}{Player-one button bits. Player two uses the same layout shifted
into higher bits (\code{PRG32\_P2\_*}).}
\label{tab:buttons}
\end{center}

The reason a bitmask is used, rather than seven separate variables, is that it
mirrors how the hardware actually presents the buttons: as bits in a register. A
student who understands the input bitmask understands, in miniature, how a
processor reads any memory-mapped peripheral. We will use this fact in both
tracks.

\begin{prgwarn}
Test a button with bitwise AND (\code{\&}), not equality (\code{==}). The mask may
hold several pressed buttons at once, so \code{input \& PRG32\_BTN\_LEFT} asks
``is the left bit set, regardless of the others?''---which is what you want.
Writing \code{input == PRG32\_BTN\_LEFT} would only be true when \emph{exactly}
left and nothing else is pressed.
\end{prgwarn}

\section{Drawing: the graphics calls you will use most}

The framework offers a small drawing interface, and a beginner needs only a
handful of its functions. Each of the calls below works identically in assembly
(passing arguments in registers) and in C (passing arguments normally), and each
works identically on QEMU and on the ESP32\nobreakdash-C6 \citep{prg32qemu}.

\begin{center}
\begin{tabularx}{\linewidth}{@{}lX@{}}
\toprule
Call & What it does \\
\midrule
\code{prg32\_gfx\_clear(color)} & fill the whole viewport with one colour \\
\code{prg32\_gfx\_pixel(x,y,color)} & set a single pixel \\
\code{prg32\_gfx\_rect(x,y,w,h,color)} & draw a filled rectangle \\
\code{prg32\_gfx\_text8(x,y,s,fg,bg)} & draw a string of text \\
\code{prg32\_gfx\_present()} & show the finished frame on the display \\
\code{prg32\_sprite\_hitbox(\dots)} & test whether two rectangles overlap \\
\bottomrule
\end{tabularx}
\captionof{table}{The core drawing and collision calls. The full graphics
interface is catalogued in Appendix~\ref{app:api}.}
\label{tab:gfxcore}
\end{center}

The rectangle call, \code{prg32\_gfx\_rect}, is the workhorse of early examples:
a paddle is a wide rectangle, a ball is a small square, a wall is a tall thin
rectangle. The collision call, \code{prg32\_sprite\_hitbox}, takes the position
and size of two rectangles and returns a non-zero value when they overlap, which
is exactly the test a beginner needs for ``did the ball hit the paddle?''

\section{Sound: a beep when something happens}

Audio in \prg starts as simply as it can: \code{prg32\_audio\_beep(hz, ms)} plays
a tone of a given frequency for a given number of milliseconds \citep{prg32abi}.
On the physical board this drives a small buzzer; on QEMU the call is accepted but
silent, because the emulator does not model the buzzer. The framework's own
advice, which this book repeats, is to play a beep only on an \emph{event}---a
collision, a score, a button press---and never on every frame, because a tone
restarted sixty times a second is an unpleasant buzz rather than a sound effect.

\begin{prgwarn}
Do not call \code{prg32\_audio\_beep} every frame. Trigger sound on events only.
A beep in \code{update} that fires whenever the ball is in flight will play
continuously; a beep that fires only in the frame where a collision is newly
detected will sound like a game.
\end{prgwarn}

\section{Cartridges: how a game reaches the machine}

A finished \prg game is packaged as a \emph{cartridge}: a file ending in
\fname{.prg32} that contains your compiled \riscv code together with a small
header \citep{prg32cartfmt}. The resident firmware---flashed onto the board once,
at the start of the course---knows how to load a cartridge and call its three
functions. This separation is what lets you iterate quickly: you flash the
firmware a single time, then upload cartridge after cartridge without reflashing,
either over the board's own Wi-Fi network or, on the emulator, by staging the
cartridge into the QEMU flash image.

\begin{center}
\begin{tikzpicture}[font=\sffamily\scriptsize,node distance=4mm,
  s/.style={rounded corners=2pt,draw=prgsteel,fill=prgpanel,minimum height=10mm,minimum width=24mm,align=center},
  ar/.style={-{Stealth[length=2mm]},draw=prgink!70}]
  \node[s] (src) {\code{game.S}\\ or \code{game.c}};
  \node[s,right=8mm of src] (tc) {\riscv\\toolchain};
  \node[s,right=8mm of tc] (cart) {\code{.prg32}\\cartridge};
  \node[s,right=8mm of cart] (rt) {\prg\\runtime};
  \node[s,right=8mm of rt] (loop) {init / update\\/ draw};
  \draw[ar] (src) -- (tc);
  \draw[ar] (tc) -- (cart);
  \draw[ar] (cart) -- (rt);
  \draw[ar] (rt) -- (loop);
\end{tikzpicture}
\captionof{figure}{From source to running game. The same pipeline serves
assembly and C, hardware and emulator \citep{prg32repo}.}
\label{fig:pipeline}
\end{center}

The cartridge header begins with the magic bytes \code{PRG2} and records, among
other things, the offsets of the \code{init}, \code{update}, and \code{draw}
functions, so the firmware knows where to jump \citep{prg32cartfmt}. You will not
write this header by hand; a tool called \fname{prg32\_game.py} builds it for you,
and Appendix~\ref{app:api} documents the tool's command line. For now the
important idea is conceptual: a cartridge is your three functions, compiled to
native code, wrapped in just enough metadata for the firmware to run them.

\section{The two targets, concretely}

You will build every game for two targets. The commands differ only in a build
directory and a configuration profile; the source does not change at all
\citep{prg32qemu}.

For the \textbf{physical ESP32\nobreakdash-C6 board}, the build uses a directory named
\fname{build-esp32c6} and the board configuration profile:

\begin{lstlisting}[style=bash,caption={Building and flashing for the physical board (abbreviated; see Appendix~\ref{app:env}).}]
idf.py -B build-esp32c6 -D SDKCONFIG_DEFAULTS=sdkconfig.defaults set-target esp32c6
idf.py -B build-esp32c6 -D SDKCONFIG_DEFAULTS=sdkconfig.defaults flash monitor
\end{lstlisting}

For the \textbf{QEMU emulator}, the build uses \fname{build-qemu} and the QEMU
profile, and targets the ESP32\nobreakdash-C3 emulator:

\begin{lstlisting}[style=bash,caption={Building and running for the desktop emulator.}]
idf.py -B build-qemu -D SDKCONFIG_DEFAULTS=sdkconfig.defaults.qemu set-target esp32c3
idf.py -B build-qemu -D SDKCONFIG_DEFAULTS=sdkconfig.defaults.qemu qemu --graphics monitor
\end{lstlisting}

\begin{prgwarn}
Keep the two build directories strictly separate. Flashing a \fname{build-qemu}
image to a real board produces a black screen, because the emulator build uses a
virtual display panel that does not exist on real hardware. If your board shows
nothing, the first thing to check is that you built in \fname{build-esp32c6}, not
\fname{build-qemu} \citep{prg32qemu}.
\end{prgwarn}

\begin{prgaside}
Why does the emulator target an ESP32\nobreakdash-C3 while the board is an ESP32\nobreakdash-C6? Both
are \riscv microcontrollers from the same family, and both run the same 32-bit
\riscv instructions. Espressif maintains a mature QEMU graphics path for the
ESP32\nobreakdash-C3, so \prg uses it for desktop emulation, while the physical kits are
built around the newer ESP32\nobreakdash-C6 \citep{prg32qemu}. From your game's point of
view the difference is invisible: the calling convention and the framework
interface are identical.
\end{prgaside}

\section{Design principles for first games}

The \prg learning path distils years of classroom experience into a few rules,
and they are worth stating plainly because beginners who follow them succeed and
beginners who ignore them struggle \citep{prg32repo,prg32tutorial}.

The first rule is \emph{start with returns that do nothing}. Write \code{init},
\code{update}, and \code{draw} so that each simply returns immediately, build it,
and confirm the framework runs without crashing. You now have a skeleton that
compiles and runs; everything after is addition.

The second rule is \emph{add one thing at a time}. One variable for position.
Then draw one object at that position. Then read input and change the position.
Then add one boundary check. Then one collision. Each step is small enough to
verify by eye before the next begins.

The third rule is \emph{sound comes last}. Get the core loop working and visible
before adding audio, because a beep that fires at the wrong time is a symptom, not
a feature, and you want the visual logic solid first.

\begin{prgconcept}
The intended rhythm of work is small and visible: change one thing, run it,
observe the result, and write down what changed \citep{prg32repo}. This is not
merely advice for beginners; it is how experienced systems programmers actually
work when the machine is unforgiving.
\end{prgconcept}

\section{A pedagogical map of the example games}

The platform ships with a family of complete example games, each provided in
three forms: an ASCII assembly version, a graphics assembly version, and a C
version \citep{prg32repo}. They are arranged, roughly, from simplest to richest,
and this book draws its worked examples from them in that order.

\begin{center}
\begin{tabularx}{\linewidth}{@{}lX@{}}
\toprule
Example & What it teaches first \\
\midrule
\code{pong} & state, input, boundaries, one collision \\
\code{breakout} & arrays of objects, removing objects on hit \\
\code{space\_invaders} & many moving objects, simple enemy waves \\
\code{pacman} & grid movement, a tile map \\
\code{asteroids} & wrap-around motion, rotation \\
\code{tetris} & falling pieces, a fixed grid, line clears \\
\code{platformer} & tile flags, gravity, camera follow \\
\code{raycaster} & fixed-point arithmetic, a pseudo-3D corridor \\
\code{wing\_commander} & two layered playfields, a cockpit \\
\bottomrule
\end{tabularx}
\captionof{table}{The example games, ordered by the concepts they introduce. This
book builds Pong from scratch in both tracks and uses the others as references
\citep{prg32repo,prg32teaching}.}
\label{tab:games}
\end{center}

We will build Pong together, line by line, in both the assembly track
(Chapters~\ref{ch:asm1}--\ref{ch:asmcap}) and the C track
(Chapters~\ref{ch:c1}--\ref{ch:ccap}), because it is the smallest game that still
contains every essential idea. The richer examples then become readable: once you
have written Pong, the Breakout source is Pong with an array of bricks, and the
platformer source is Pong with gravity and a tile map.

\begin{prgcheck}
Before starting a track, make sure you can answer these without looking back:
What are the three functions every \prg game provides? In which direction does
$y$ grow on the screen? How do you test whether one button is pressed? What is a
cartridge? Why must the QEMU and board builds use separate directories?
\end{prgcheck}
